% GENERATED FILE. Edit canonical lessons, then run npm run generate:books.
% canonical-source-hash: fnv1a64:040b4a9eb03ace92
% canonical-lessons: LA-C28-dies

\chapter{Day}
\label{ch:la-day}

This chapter is generated from the canonical micro-lessons used by Language Ladder.
Each section stays independently resumable and preserves the authored prerequisite order.

\section[diēs]{diēs — the word you already know, now standing alone}

\label{lesson:LA-C28-dies}



\begin{quote}
\textbf{Warm-up.} {[}PAUSE 2s{]} You've used this word twice already without a lesson of its own — inside \emph{diēs Lūnae} and \emph{medius diēs}. Time to meet it directly.
\end{quote}

\subsection*{You'll want to know: diēs — "day"}
\textbf{Diēs} — "\textbf{day}" — is a genuinely common, everyday Classical Latin word, not a rare or bookish one. You've already met it embedded in:

\begin{itemize}
\raggedright
  \item \textbf{diēs Lūnae} ("Monday," literally "day of the Moon" — Chapter 5)
  \item \textbf{medius diēs} ("noon," literally "mid-day" — Chapter 12)
\end{itemize}

Now it's the headword itself.

\begin{cousinweb}
Diēs has its own derivative, \textbf{diurnum} ("daytime, daily"), which gives English \textbf{diurnal} directly — and, through French's \emph{jour}, \emph{journée}, and \emph{jornel} (all descendants of \emph{diurnum} in their own right), \textbf{journal} and \textbf{journey} as well (a "journey" was originally "a \textbf{day's} travel"). \textbf{Diary} takes a \textbf{separate} path: it comes from a \textbf{sibling} derivative of \emph{diēs}, \textbf{diārium} ("daily allowance"), not from \emph{diurnum} itself — two related but distinct branches off the same root. Further out, a genuine \textbf{Proto-Indo-European} root, \emph{\textbf{*dyēws-}} ("to shine, sky"), connects \emph{diēs} to Latin's own \textbf{deus} ("god") and to \textbf{Iuppiter} (Jupiter, literally "sky-father") — the same root that shines through the sky-god imagery of gods across the Indo-European world.
\end{cousinweb}

\subsection*{Guided Practice}
{[}PAUSE 1s{]}

\begin{itemize}
\raggedright
  \item {[}YOU SAY: "diēs" — "day"{]}
  \item {[}YOU SAY: where you've already met it — diēs Lūnae, medius diēs{]}
  \item {[}YOU SAY: the English cousins — diurnal, diary, journal, journey{]}
\end{itemize}

\begin{tcolorbox}[breakable,colback=teal!4,colframe=teal!35!black,title={Wrap-up recall}]
{[}PAUSE 3s{]} Where have you already met \textbf{diēs} twice before this lesson? (\textbf{Diēs Lūnae}, "Monday"; \textbf{medius diēs}, "noon.") What English word comes from diēs's derivative \emph{diurnum} directly? (\textbf{Diurnal}.) What two words come through French descendants of \emph{diurnum} — \emph{jour}, \emph{journée}, \emph{jornel}? (\textbf{Journal, journey}.) Does \textbf{diary} also come from \emph{diurnum}? (\textbf{No} — a separate sibling derivative, \emph{diārium}, "daily allowance.")
\end{tcolorbox}
